\begin{theorem}[Integral torsor classes and the parameter period]\hypertarget{tag-02YR}{}\label{night:P5.6}\label{stmt:P5.6}
\StatementRecord{night:P5.6}{proved}{02YR}{the definitions and constructions stated in the proof}

For fixed admissible twist data, $X_c\simeq X_{c'}$ precisely when
$c-c'\in2\mathbb Z$.  This fixes the integral period of $c$; the
construction still permits an additive choice of $\beta^{\rm part}$.
\end{theorem}
\begin{proof}
\begin{enumerate}[label=\arabic*.,leftmargin=*]
\item The analytically computed complete linear system
$|\omega_{X_c}^{-2}|$ is the marked map $f_c$, by
\cref{wb:thm:canonical-ring-recovery}. A biholomorphism preserves its
canonical bundle, the complete space of sections, and its evaluation
map. Consequently it induces $\alpha\in\operatorname{PGL}_2(\mathbb C)$
with $f_{c'}F=\alpha f_c$. The three special fibres have distinct
intrinsic types: the cusp is reduced and nonnormal, while the other
two have multiplicities three and four. All remaining fibres are
smooth reduced tori. Thus $\alpha$ fixes $p_0,p_1,p_2$ individually.
In the marked coordinate, a fractional-linear map fixing
$\infty,0,1$ has zero lower-left and upper-right entries and equal
diagonal entries; hence $\alpha=\mathrm{id}_B$.

The centralizer computation of \cref{wb:pshift:lem:centralizer}
then gives $P=aI+bE_{41}$ with $a=\pm1$ and $b=a(c-c')$.
The root-fibre character calculation in that proposition excludes
$a=-1$. Thus $b=c-c'$ is an integer.

\item Compute $H^2(\mathcal B,\mathcal H)$ with a pair of pants, two root discs,
and a cusp disc.  The pair of pants has cochains $\Lambda\to\Lambda^2$
with differential $((A_1-I),(A_2-I))$.  A root disc has the periodic
cyclic resolution with alternating differentials $A_j-I,N_j$; the
cusp disc has $\ker(M_0-I)$ in degree zero.  An annulus has
$\Lambda\xrightarrow{A_j-I}\Lambda$.  The degree-one restriction from
the pair of pants is
\[
 (x_1,x_2)\longmapsto(x_1,x_2,-M_0(x_1+A_1x_2)).
\]
The total-complex ranks in degrees zero through three are $14,28,20,8$.
The resulting integer differentials have sizes
$28\times14$, $20\times28$ and $8\times20$, and their consecutive
products vanish by the restriction identities.
Smith decompositions with unimodular changes of basis identify each
integral kernel before quotienting by the preceding image, giving
$H^0=\mathbb Z$, $H^1=0$, $H^2=\mathbb Z\oplus\mathbb Z/2$.


\item The degree-two presentation also reduces by hand.  The invariant
lattices have bases $(\widehat\delta,\varepsilon_j)$, where
$\varepsilon_1=(1,2,-4,0)$ and $\varepsilon_2=(1,3,-3,0)$.
The five generators $(\delta_1,\varepsilon_1,\delta_2,\varepsilon_2,h_c)$
have relations
\begin{equation}\hypertarget{tag-01OF}{}\label{eq:P5H2}
 \delta_1=0,\qquad 2\delta_2=0,\qquad
 3\varepsilon_1=h_c=4\varepsilon_2.
\end{equation}
For the constant sheaf $\mathbb Z$, the same computation has generators
$z_1,z_2,h_c$ and relations $3z_1=h_c=4z_2$; its primitive generator is
$z_1-z_2$.  Inclusion $\mathbb Z\widehat\delta\to\mathcal H$ sends it
to the nonzero order-two class $\delta_2$.  The character $\gamma$
sends $\varepsilon_j$ to $z_j$ and kills $\delta_j$.

\item Both $H^1(\mathcal B,E)$ and $H^2(\mathcal B,E)$ vanish, by
$\rho_*E=\mathcal O\oplus\mathcal O(-1)$ and the dimension of the coarse
curve.  Consequently \eqref{eq:P5exponential} identifies compactified
translation torsors $H^1(\mathcal B,\mathcal A)$ with
$H^2(\mathcal B,\mathcal H)$.  This identification takes logarithms of
local translation transitions and their lattice-valued \v Cech
coboundaries; the two vanishings also reverse that construction.

\item The cyclic norm of the translation $v_j/m_j$ is $v_j$, while the cusp
collar has the opposite boundary orientation.  Thus the torsor
character is
\[
 l_1z_1+l_2z_2-l_0h_c
 =(4l_1+3l_2-12l_0)(z_1-z_2)=-p(z_1-z_2).
\]
The shear $S_b-I=b\widehat\delta\otimes\gamma$ changes its class by
$-bp\delta_2$.  Freeness forces $l_2$ odd and thus $p$ odd.  The class
is preserved precisely for even $b$.

\item For an integral shear the Jacobian charts themselves identify: the
period lattices differ by $S_b$, and the bounded cusp matrix changes by
an integer in its $\beta$ entry.  Every toric unit
$\exp(2\pi iC\lambda)$ is therefore unchanged.  The fan, lattice, and
toric gluings are the same, and on root discs the basis change
identifies the smooth Jacobian and its cyclic action.  Equality of the
torsor classes now gives a holomorphic isomorphism for the fixed sheaf
$\mathcal A$, which consists precisely of translations extending over
the cusp.  It extends on all three special charts.  Conversely every
isomorphism preserves the class, so forces $b$ even.
\end{enumerate}
\end{proof}